\documentclass[11pt]{article}

\usepackage[margin=1.05in]{geometry}
\usepackage{amsmath,amssymb,amsthm,mathtools}
\usepackage{booktabs}
\usepackage{enumitem}
\usepackage[hidelinks]{hyperref}

\newtheorem{theorem}{Theorem}[section]
\newtheorem{proposition}[theorem]{Proposition}
\newtheorem{lemma}[theorem]{Lemma}
\newtheorem{corollary}[theorem]{Corollary}
\theoremstyle{definition}
\newtheorem{definition}[theorem]{Definition}
\theoremstyle{remark}
\newtheorem{remark}[theorem]{Remark}

\newcommand{\E}{\mathbb E}
\newcommand{\Pp}{\mathbb P}
\newcommand{\one}{\mathbf 1}
\newcommand{\cP}{\mathcal P}
\newcommand{\cM}{\mathcal M}
\newcommand{\cB}{\mathcal B}
\newcommand{\cD}{\mathcal D}
\newcommand{\cZ}{\mathcal Z}
\newcommand{\RC}{\phi}

\title{Sharp Boolean--Tutte Completion for\\Finite Random-Cluster Gate Inequalities}
\author{}
\date{September 4, 2026}

\begin{document}
\maketitle

\begin{abstract}
The previously proposed Christoffel--Tutte hierarchy applies
Cauchy--Schwarz to intersection indicators.  It is valid, but it is not the
strongest conclusion carried by the same intersection moments.  We replace
it by an exact finite linear program on gate signatures.  The new bound is
the largest lower bound valid for every finite event system with the stated
moments, dominates the Christoffel/Rayleigh bound at every level, and is
strictly stronger in general.  Its permutation-invariant quotient is a
finite factorial-moment program.  At order two this gives the sharp
Dawson--Sankoff bound, which strictly improves the second-moment bound and
is already exact for two gates.  For finite random-cluster measures the
same construction has an order-free component-partition expansion.  Every
dual certificate then gives a coefficientwise nonnegative Tutte-polynomial
inequality, valid simultaneously for every cluster weight $q>0$.

The requested conclusion for all vertex-transitive graphs is not asserted.
Taken literally it is false when $p_c=1$ (for example on $\mathbb Z$), and
the corrected assertion for all transitive graphs with $p_c<1$ remains
strictly stronger than the results proved here.  The Hutchcroft/Kruskal
program supplied with the question still requires its stated
first-reconvergence or nested anti-congestion estimate.  The finite theorem
below does not imply that estimate without new quantitative input.
\end{abstract}

\section{Why the Gram hierarchy is not optimal}

Let $(\Omega,\Pp)$ be a finite probability space, let $I$ be a nonempty
finite set, and let $(E_i)_{i\in I}$ be events.  Put
\[
 D=\bigcup_{i\in I}E_i.
\]
Suppose that $X\geq0$ and $X=0$ off $D$.  For every nonempty $S\subseteq I$
write
\[
 J_S=\one_{\cap_{i\in S}E_i},\qquad
 m_S=\E[XJ_S].
\]
The Christoffel construction forms the matrix $M_{S,T}=m_{S\cup T}$ and
uses
\[
 \E[X]\geq \frac{(c^{\mathsf T}m)^2}{c^{\mathsf T}Mc}.
 \tag{1.1}
\]
This is Cauchy--Schwarz and is correct.  It nevertheless discards the most
important fact about the indicators $J_S$: they are monomials of one
Boolean signature.  Retaining that fact gives an exact linear program.

The loss is already visible for two events.  If
\[
 a=\E[X\one_{E_1}],\qquad d=\E[X\one_{E_2}],\qquad
 c=\E[X\one_{E_1\cap E_2}],
\]
then the three numbers used in the degree-one Gram matrix determine the
answer exactly:
\[
 \E[X]=a+d-c.
 \tag{1.2}
\]
By contrast, when $ad>c^2$, the optimized Gram value is
\[
 \frac{ad(a+d-2c)}{ad-c^2},
\]
and its deficit from (1.2) is
\[
 \frac{c(a-c)(d-c)}{ad-c^2}.
 \tag{1.3}
\]
Thus a genuine strengthening must use the Boolean constraints rather than
another reformulation of the same positive-semidefinite matrix.

\section{The sharp signature program}

For $\omega\in\Omega$, define its gate signature
\[
 \sigma(\omega)=\{i\in I:\omega\in E_i\}.
\]
For every nonempty $\sigma\subseteq I$, put
\[
 z_\sigma=\E\bigl[X;\sigma(\omega)=\sigma\bigr].
 \tag{2.1}
\]
Then
\[
 \E[X]=\sum_{\varnothing\ne\sigma\subseteq I}z_\sigma,
 \qquad
 m_S=\sum_{\sigma\supseteq S}z_\sigma.
 \tag{2.2}
\]

Fix a family
\(
 \cM\subseteq\cP(I)\setminus\{\varnothing\}
\)
of moments that are to be retained.  For a vector
$m=(m_S)_{S\in\cM}$ define
\[
 \cP_{\cM}(m)=
 \left\{(u_\sigma)_{\varnothing\ne\sigma\subseteq I}:
 u_\sigma\geq0,
 \ \sum_{\sigma\supseteq S}u_\sigma=m_S
 \text{ for every }S\in\cM
 \right\}
\]
and
\[
 \cB_{\cM}(m)=
 \min_{u\in\cP_{\cM}(m)}
 \sum_{\varnothing\ne\sigma\subseteq I}u_\sigma.
 \tag{2.3}
\]
The minimum exists whenever the data arise from an event system: the
actual vector (2.1) is feasible, and intersecting the feasible set with the
sublevel set determined by its total mass makes the optimization compact.

\begin{theorem}[Sharp Boolean moment completion]\label{thm:signature}
For every finite event system and every $X\geq0$ supported on $D$,
\[
 \boxed{\E[X]\geq\cB_{\cM}(m).}
 \tag{2.4}
\]
Moreover,
\[
 \boxed{
 \cB_{\cM}(m)=
 \max\left\{
  \sum_{S\in\cM}c_Sm_S:
  \sum_{\substack{S\in\cM\\S\subseteq\sigma}}c_S\leq1
  \text{ for every }\varnothing\ne\sigma\subseteq I
 \right\}.}
 \tag{2.5}
\]
The maximum is attained.  Among bounds depending only on the listed
moments and valid for all finite event systems, (2.4) is best possible.
\end{theorem}

\begin{proof}
The actual signature masses $z_\sigma$ satisfy (2.2), so they are feasible
in (2.3).  Their objective is $\E[X]$, proving (2.4).

Let $A$ be the zero-one matrix with rows indexed by $S\in\cM$, columns
indexed by nonempty $\sigma\subseteq I$, and entries
$A_{S,\sigma}=\one_{\{S\subseteq\sigma\}}$.  The primal program is
\[
 \min\{\mathbf1^{\mathsf T}u:Au=m,\ u\geq0\}.
\]
Finite-dimensional linear-programming duality gives
\[
 \max\{c^{\mathsf T}m:A^{\mathsf T}c\leq\mathbf1\},
\]
which is exactly (2.5).  For completeness, weak duality follows from
\[
 c^{\mathsf T}m=c^{\mathsf T}Au
 =(A^{\mathsf T}c)^{\mathsf T}u
 \leq\mathbf1^{\mathsf T}u.
\]
For the reverse inequality, let $\beta$ be the primal minimum and consider
the finitely generated closed cone
\[
 K=\{(Au,\mathbf1^{\mathsf T}u+t):u\geq0,\ t\geq0\}.
\]
For every $\varepsilon>0$, the point $(m,\beta-\varepsilon)$ is outside
$K$.  Separating this point from $K$ gives $(y,s)$ with $s>0$,
$A^{\mathsf T}y+s\mathbf1\geq0$, and
$y^{\mathsf T}m+s(\beta-\varepsilon)<0$.  The vector
$c=-y/s$ is dual feasible and satisfies
$c^{\mathsf T}m>\beta-\varepsilon$.  Hence the dual supremum is $\beta$;
the standard finite-polyhedron attainment argument gives a maximizer.

It remains to justify the optimality claim.  Let $u$ be any feasible
vector in (2.3).  Make one atom $\omega_\sigma$ for every nonempty
signature, together with one dummy atom, give all these atoms equal
positive probability, put $\omega_\sigma\in E_i$ exactly when
$i\in\sigma$, and choose the value of $X(\omega_\sigma)$ so that its
expected mass is $u_\sigma$; put $X=0$ at the dummy atom.  This realizes
the prescribed moment vector and has
$\E[X]=\sum_\sigma u_\sigma$.  In particular, an optimizer in (2.3) is
realized by a finite event system.  No universally valid function of the
listed moments can therefore give a larger lower bound.
\end{proof}

Adding moments can only shrink the feasible set in (2.3).  Thus the bounds
are monotone under enrichment of $\cM$.  If all nonempty subsets of $I$ are
included, M\"obius inversion determines every $z_\sigma$, and
\[
 \boxed{
 \cB_{\cP(I)\setminus\{\varnothing\}}(m)
 =\sum_{\varnothing\ne S\subseteq I}(-1)^{|S|+1}m_S
 =\E[X].}
 \tag{2.6}
\]
This is finite inclusion--exclusion, now identified as the terminal point
of an information-optimal hierarchy.

\subsection{Using bounded observables}

When $0\leq X\leq1$, one can retain the gate-signature probabilities
\[
 h_\sigma=\Pp(\sigma(\omega)=\sigma).
\]
Then $0\leq z_\sigma\leq h_\sigma$.  Define the capped program
\[
 \cB^{\mathrm{cap}}_{\cM}(m,h)=
 \min\left\{
  \sum_\sigma u_\sigma:Au=m,\ 0\leq u_\sigma\leq h_\sigma
 \right\}.
 \tag{2.7}
\]

\begin{proposition}[Capped dual]\label{prop:capped}
If $0\leq X\leq1$ and $X=0$ off $D$, then
\[
 \E[X]\geq\cB^{\mathrm{cap}}_{\cM}(m,h)
 \geq\cB_{\cM}(m),
 \tag{2.8}
\]
and
\[
 \boxed{
 \cB^{\mathrm{cap}}_{\cM}(m,h)
 =\max_{c\in\mathbb R^{\cM}}
 \left\{
 c^{\mathsf T}m-
 \sum_{\sigma\ne\varnothing}h_\sigma
 \bigl((A^{\mathsf T}c)_\sigma-1\bigr)_+
 \right\}.}
 \tag{2.9}
\]
This is the best universal lower bound determined by $(m,h)$.
\end{proposition}

\begin{proof}
The actual $z$ is feasible in (2.7), and the capped feasible set is
contained in the uncapped one.  For a multiplier $c$ for $Au=m$, minimizing
the Lagrangian over the box $0\leq u\leq h$ gives
\[
 c^{\mathsf T}m+
 \sum_\sigma\min_{0\leq u_\sigma\leq h_\sigma}
 \bigl(1-(A^{\mathsf T}c)_\sigma\bigr)u_\sigma,
\]
which is the expression in (2.9).  Finite linear-programming duality gives
equality.  Conversely, every vector $0\leq u\leq h$ can be realized on a
space with signature masses $h_\sigma$ by setting
$X=u_\sigma/h_\sigma$ on the $\sigma$-atom (and setting it to zero when
$h_\sigma=0$).  This proves optimality.
\end{proof}

\section{Comparison with the Christoffel bound}

\begin{theorem}[The Boolean program dominates the Gram program]
Let $\mathcal F$ be a family of nonempty subsets of $I$, and suppose
$\cM$ contains every set $S\cup T$ with $S,T\in\mathcal F$.  Form
\[
 m^{\mathcal F}_S=m_S,
 \qquad M^{\mathcal F}_{S,T}=m_{S\cup T}.
\]
Then
\[
 \boxed{
 \cB_{\cM}(m)\geq
 \sup_{c^{\mathsf T}M^{\mathcal F}c>0}
 \frac{(c^{\mathsf T}m^{\mathcal F})^2}
      {c^{\mathsf T}M^{\mathcal F}c}.}
 \tag{3.1}
\]
The inequality can be strict.
\end{theorem}

\begin{proof}
Take any feasible signature vector $u$ in (2.3), and regard $u$ as a
finite measure on the nonempty signatures.  On this space let
$J_S(\sigma)=\one_{\{S\subseteq\sigma\}}$.  Its listed first and second
moments are exactly $m^{\mathcal F}$ and $M^{\mathcal F}$.  Applying
Cauchy--Schwarz under this measure shows
\[
 \sum_\sigma u_\sigma\geq
 \frac{(c^{\mathsf T}m^{\mathcal F})^2}
      {c^{\mathsf T}M^{\mathcal F}c}
\]
for every admissible $c$.  Minimize the left side over feasible $u$ and
then maximize the right side over $c$.  Strictness follows from the
two-event calculation (1.2)--(1.3).
\end{proof}

For the four-cycle data quoted in the preceding draft,
\[
 a=d=\frac{325}{398},\qquad c=\frac{625}{796}.
\]
The Gram value is
\[
 \frac{8450}{10149}\approx0.832594,
\]
whereas the Boolean program, using exactly the same three moments, gives
\[
 a+d-c=\frac{675}{796}\approx0.847990.
\]
Thus it recovers the exact answer one full level earlier than the proposed
Gram hierarchy.

\section{The sharp gate-count hierarchy}

The full signature program distinguishes the gates.  A much smaller
program results if only permutation-invariant moment sums are retained.
Let $n=|I|$ and
\[
 N=\sum_{i\in I}\one_{E_i}.
\]
On the support of $X$, one has $1\leq N\leq n$.  Define the factorial
moments
\[
 S_j=\E\left[X\binom Nj\right]
 =\sum_{\substack{T\subseteq I\\|T|=j}}m_T,
 \qquad 1\leq j\leq n.
 \tag{4.1}
\]
For $1\leq d\leq n$, put
\[
 \cD_d(S)=
 \min\left\{
  \sum_{k=1}^n u_k:
  u_k\geq0,\quad
  \sum_{k=1}^n\binom kj u_k=S_j\ (1\leq j\leq d)
 \right\}.
 \tag{4.2}
\]

\begin{theorem}[Sharp factorial-moment completion]\label{thm:count}
For every $d$,
\[
 \boxed{\E[X]\geq\cD_d(S),}
 \tag{4.3}
\]
where
\[
 \boxed{
 \cD_d(S)=
 \max\left\{
  \sum_{j=1}^d a_jS_j:
  \sum_{j=1}^d a_j\binom kj\leq1
  \text{ for }k=1,\ldots,n
 \right\}.}
 \tag{4.4}
\]
This is the best bound based only on $S_1,\ldots,S_d$.  The sequence is
nondecreasing in $d$, and
\[
 \boxed{\cD_n(S)=\E[X].}
 \tag{4.5}
\]
\end{theorem}

\begin{proof}
Put
\[
 u_k=\E[X;N=k].
\]
These nonnegative numbers are feasible in (4.2), their total is $\E[X]$,
and (4.4) is the finite LP dual.  Every feasible $(u_k)$ is realized by a
finite event system supported on signatures of the indicated sizes, so the
bound is information-optimal.  Enriching the moment list shrinks the
primal feasible set.  Finally,
\[
 \sum_{j=1}^n(-1)^{j+1}\binom kj=1
 \qquad(1\leq k\leq n).
\]
Using these coefficients in (4.4) gives $\E[X]$, proving (4.5).
\end{proof}

The order-two program has a closed form.  This is the classical sharp
Dawson--Sankoff lower bound.

\begin{theorem}[Sharp order-two formula]\label{thm:DS}
Assume $S_1>0$ and put
\[
 x=1+\frac{2S_2}{S_1},\qquad r=\lfloor x\rfloor.
 \tag{4.6}
\]
Then $1\leq r\leq n$ and
\[
 \boxed{
 \cD_2(S)=
 \frac{2S_1}{r+1}-\frac{2S_2}{r(r+1)}.}
 \tag{4.7}
\]
If $S_1=0$, then $X=0$ almost surely and all quantities vanish.  The bound
(4.7) dominates the Cauchy--Schwarz estimate
\[
 \frac{S_1^2}{S_1+2S_2},
 \tag{4.8}
\]
and their difference is
\[
 \boxed{
 \frac{S_1(x-r)(r+1-x)}{xr(r+1)}\geq0.}
 \tag{4.9}
\]
It is strict unless $x$ is an integer.
\end{theorem}

\begin{proof}
For every integer $k\geq1$ and every positive integer $r$,
\[
 1-\left(\frac{2k}{r+1}
 -\frac{2\binom{k}{2}}{r(r+1)}\right)
 =\frac{(k-r)(k-r-1)}{r(r+1)}\geq0.
 \tag{4.10}
\]
The last inequality holds because no integer lies strictly between $r$ and
$r+1$.  Multiplying (4.10) by $u_k$ and summing proves that the right side
of (4.7) is a valid dual bound.

To prove optimality, let $\eta=x-r$.  Since $x$ is the mean of $k$ under
the probability weights $ku_k/S_1$, one has $1\leq x\leq n$ and
$0\leq\eta<1$, except that $x=n$ simply gives $\eta=0$ and $r=n$.
Define
\[
 u_r=\frac{S_1(1-\eta)}r,
 \qquad
 u_{r+1}=\frac{S_1\eta}{r+1},
 \tag{4.11}
\]
with the second term omitted when $r=n$, and set all other $u_k$ to zero.
Then
\[
 \sum k u_k=S_1,
 \qquad
 \sum\binom{k}{2}u_k=S_2.
\]
Its total mass is exactly the right side of (4.7), so the primal and dual
values agree.  Since $S_1+2S_2=xS_1$, direct subtraction of (4.8) from
(4.7) gives (4.9).
\end{proof}

When $n=2$, write $a=m_{\{1\}}$, $d=m_{\{2\}}$, and
$c=m_{\{1,2\}}$.  Then $S_1=a+d$ and $S_2=c$; Theorem~\ref{thm:DS}
gives $a+d-c=\E[X]$.  Thus the order-two integer constraint closes exactly
the deficit (1.3) left by the continuous Cauchy--Schwarz relaxation.

\section{Application to the random-cluster model}

Let $G=(V,E)$ be a finite loopless multigraph with edge parameters
$p_e\in[0,1]$ and cluster weight $q>0$.  Its random-cluster law is
\[
 \RC_{G,p,q}(\omega)=\frac1{Z_{G,p,q}}
 q^{k_G(\omega)}
 \prod_{e\in\omega}p_e
 \prod_{e\notin\omega}(1-p_e).
 \tag{5.1}
\]
No correlation assumption and no restriction on $q$ is needed below.

Fix nonempty vertex sets $R,A\subseteq V$, and let
\[
 C_R=\bigcup_{r\in R}C_r,
 \qquad D=\{R\leftrightarrow A\}.
\]
For each $a\in A$, take the single gate event
\[
 E_a=\{R\leftrightarrow a\}.
 \tag{5.2}
\]
Then $D=\bigcup_{a\in A}E_a$.  In particular, only $|A|$ Boolean gates
are required; indexing instead by $R\times A$ is unnecessary.

Let $F:2^V\to\mathbb R$ be arbitrary and put
\[
 m_F=\min_{K\subseteq V}F(K),\qquad
 X_F=(F(C_R)-m_F)\one_D.
 \tag{5.3}
\]
For nonempty $S\subseteq A$, define
\[
 m_S(F)=\RC\left[F(C_R)-m_F;
 R\leftrightarrow a\text{ for every }a\in S\right].
 \tag{5.4}
\]

\begin{theorem}[Sharp random-cluster gate completion]\label{thm:RC}
For every $q>0$, every moment family $\cM$, and every real function $F$,
\[
 \boxed{
 \RC[F(C_R);D]-m_F\RC(D)
 \geq \cB_{\cM}\bigl((m_S(F))_{S\in\cM}\bigr).}
 \tag{5.5}
\]
The capped version of Proposition~\ref{prop:capped} applies whenever
$0\leq F(C_R)-m_F\leq1$.  The factorial-moment bounds of
Theorems~\ref{thm:count} and~\ref{thm:DS} apply after summing (5.4) over
sets of equal size.  With every nonempty $S\subseteq A$ retained, (5.5) is
an equality after finitely many levels.
\end{theorem}

\begin{proof}
Apply Theorem~\ref{thm:signature} to the finite edge-configuration space,
the events (5.2), and the nonnegative variable (5.3).  Equation (5.4) is
exactly the corresponding intersection moment.  The remaining assertions
follow from Proposition~\ref{prop:capped} and
Theorems~\ref{thm:count}--\ref{thm:DS}.
\end{proof}

For a target vertex $b$, take $F(K)=\one_{\{b\in K\}}$.  Then
\[
 \alpha_b=\RC(R\leftrightarrow b,\ R\leftrightarrow A)
\]
and
\[
 p_S(b)=\RC(R\leftrightarrow b,\ R\leftrightarrow a
              \text{ for every }a\in S).
 \tag{5.6}
\]
Consequently,
\[
 \boxed{\alpha_b\geq\cB_{\cM}\bigl((p_S(b))_{S\in\cM}\bigr).}
 \tag{5.7}
\]
This bound is valid for $q<1$, $q=1$, and $q>1$.  For a single source
$R=\{o\}$, the moment in (5.4) is gate-readable: on the event in (5.4),
$C_a=C_o$ for every $a\in S$.

\section{Order-free Tutte certificates}

The random-cluster specialization has a further algebraic structure.  Use
the exact cluster-law reduction: delete edges with $p_e=0$, contract the
components formed by edges with $p_e=1$, delete loops, and replace a
parallel class $B$ by one edge of parameter
\[
 1-\prod_{e\in B}(1-p_e).
\]
It is therefore enough to consider a simple graph with $0<p_e<1$.
Put
\[
 v_e=\frac{p_e}{1-p_e}.
\]
For a nonempty vertex set $B$, define
\[
 \kappa(B)=
 \sum_{\substack{T\subseteq E(B)\\(B,T)\text{ connected}}}
 \prod_{e\in T}v_e,
 \tag{6.1}
\]
with $\kappa(\{v\})=1$.  For a partition $\pi$ of $V$, put
\[
 w(\pi)=\prod_{B\in\pi}\kappa(B),
 \qquad
 \cZ(q)=\sum_\pi q^{|\pi|}w(\pi).
 \tag{6.2}
\]
The open-component partition $\Pi$ then satisfies
\[
 \RC(\Pi=\pi)=\frac{q^{|\pi|}w(\pi)}{\cZ(q)}.
 \tag{6.3}
\]
Indeed, after the global factor $\prod_e(1-p_e)$ is removed, a
configuration with component partition $\pi$ is exactly an independent
choice of a connected spanning edge set inside every block of $\pi$.

For a partition $\pi$, let
\[
 B_R(\pi)=\bigcup\{B\in\pi:B\cap R\ne\varnothing\},
 \qquad
 \sigma_R(\pi)=A\cap B_R(\pi).
\]
For every nonempty $\sigma\subseteq A$, define the atom polynomial
\[
 W^b_\sigma(q)=
 \sum_{\substack{\pi:\ b\in B_R(\pi)\\
                         \sigma_R(\pi)=\sigma}}
 q^{|\pi|}w(\pi),
 \tag{6.4}
\]
and the gate-only atom polynomial
\[
 H_\sigma(q)=
 \sum_{\substack{\pi:\sigma_R(\pi)=\sigma}}
 q^{|\pi|}w(\pi).
 \tag{6.5}
\]
All their coefficients are nonnegative, and coefficientwise
$0\leq W^b_\sigma\leq H_\sigma$.  Put
\[
 L_b(q)=\sum_{\sigma\ne\varnothing}W^b_\sigma(q),
 \qquad
 P^b_S(q)=\sum_{\sigma\supseteq S}W^b_\sigma(q).
 \tag{6.6}
\]
Then
\[
 \alpha_b=\frac{L_b(q)}{\cZ(q)},
 \qquad
 p_S(b)=\frac{P^b_S(q)}{\cZ(q)}.
\]

\begin{theorem}[Coefficientwise Boolean--Tutte certificate]
Suppose $(c_S)_{S\in\cM}$ satisfies
\[
 a_\sigma:=
 \sum_{\substack{S\in\cM\\S\subseteq\sigma}}c_S\leq1
 \qquad(\varnothing\ne\sigma\subseteq A).
 \tag{6.7}
\]
Then
\[
 \boxed{
 L_b(q)-\sum_{S\in\cM}c_SP^b_S(q)
 =\sum_{\sigma\ne\varnothing}(1-a_\sigma)W^b_\sigma(q),}
 \tag{6.8}
\]
so the polynomial on the left has nonnegative coefficients.  In
particular, after division by $\cZ(q)$, the corresponding random-cluster
inequality holds simultaneously for every $q>0$.

More generally, for arbitrary $c$ put $u_\sigma=(a_\sigma-1)_+$.  Then
\[
 \boxed{
 L_b(q)\geq
 \sum_{S\in\cM}c_SP^b_S(q)
 -\sum_{\sigma\ne\varnothing}u_\sigma H_\sigma(q),}
 \tag{6.9}
\]
and the difference is again coefficientwise nonnegative.
\end{theorem}

\begin{proof}
Expand the right side of (6.8) using (6.6) and collect the coefficient of
each atom polynomial $W^b_\sigma$.  Condition (6.7) makes every multiplier
nonnegative.

For (6.9), the contribution of a fixed signature to the difference between
the left and right sides is
\[
 (1-a_\sigma)W^b_\sigma+u_\sigma H_\sigma.
\]
If $a_\sigma\leq1$, this is nonnegative.  If $a_\sigma>1$, it equals
\[
 (a_\sigma-1)(H_\sigma-W^b_\sigma),
\]
which is coefficientwise nonnegative.  Summing over signatures proves the
claim.
\end{proof}

This is stronger than pointwise positive semidefiniteness of the former
Christoffel--Tutte matrix in two ways: it is optimal for the chosen Boolean
moment data, and every fixed dual witness proves an identity with visibly
nonnegative coefficients rather than merely a nonnegative value at one
$q$.

\section{What this does not prove on transitive graphs}

The statement
\[
 \theta_G(p_c(G))=0
 \quad\text{for every infinite vertex-transitive graph }G
 \tag{7.1}
\]
is false without a qualification.  For nearest-neighbour bond percolation
on $\mathbb Z$, one has $p_c=1$ and $\theta(1)=1$.  The intended statement
is the conjecture obtained by adding $p_c(G)<1$.

The Hutchcroft/Kruskal notes supplied with the question do not prove that
corrected statement.  Their final displayed nested anti-congestion limit is
an additional theorem-sized hypothesis.  Calling it a consequence of
transitivity would be circular unless it is proved.

There is a second logical separation.  Theorems~\ref{thm:signature} and
\ref{thm:RC} optimally convert known finite intersection moments into a
union estimate.  They do not bound those intersection moments.  In the
replica construction, the missing first-reconvergence estimate is exactly
a bound on such overlaps.  An information-optimal union inequality cannot
create an overlap estimate that has not been supplied.

The order-two theorem does sharpen the bookkeeping in the proposed
low-congestion/high-congestion dichotomy.  If $N$ counts any finite family
of replica reconvergence events and
\[
 S_1=\E N,\qquad S_2=\E\binom N2,
\]
then Theorem~\ref{thm:DS} is the exact smallest possible value of
$\Pp(N>0)$ compatible with these two factorial moments.  Equivalently, for
every positive integer $r$,
\[
 \Pp(N>0)\geq
 \frac{2S_1}{r+1}-\frac{2S_2}{r(r+1)}.
 \tag{7.2}
\]
Thus a large first-reconvergence charge either forces a large union or a
quantitatively large pair-overlap charge.  What remains absent is a theorem
turning the latter charge, uniformly over scales, into percolation below
$p$.  The bundle-canopy example in the supplied notes shows that
unimodularity, finite mean degree, and critical survival alone do not give
that conversion.

Accordingly, the strongest rigorous conclusion is:
\begin{itemize}[leftmargin=2em]
\item the finite random-cluster hierarchy is completely strengthened by
      the sharp Boolean and capped programs above;
\item the permutation-invariant order-two relaxation has the closed sharp
      formula (4.7), and the full finite hierarchy is exact;
\item the component-partition expansion promotes every dual witness to a
      coefficientwise Tutte-polynomial inequality for all $q>0$;
\item no unconditional extension to all transitive graphs follows until
      the stated anti-congestion or an equivalent robustness theorem is
      proved.
\end{itemize}

\end{document}
